% ============================================================
% PARTE II - LA TRIADE E LA SCRITTURA A QUATTRO PARTI
% CAPITOLO 3 - LA TRIADE
% ============================================================
\part{La triade e la scrittura a quattro parti}

\chapter{La triade}
\label{chap:triadi}

Dopo la fisica del suono e gli intervalli (Capitolo \ref{chap:fondamenti}) si può costruire l'elemento base dell'armonia tonale: la \textbf{triade}. Quasi tutta la musica occidentale, dalla polifonia rinascimentale al pop, si basa sulla successione di accordi di tre note.

\section{Dalla melodia all'armonia}
\label{sec:melodia_armonia}

Finora gli intervalli sono stati letti come relazioni tra due suoni. La musica occidentale, almeno dal tardo Medioevo, oltre a far suonare le note una dopo l'altra (dimensione \textbf{orizzontale}, o melodica), sovrappone più suoni contemporaneamente (dimensione \textbf{verticale}, o armonica). Ne nascono strutture sonore con una logica propria, e l'elemento base è l'\textbf{accordo}.

\begin{definizione}[Accordo]
	\label{def:accordo}
	Un \textbf{accordo} è la sovrapposizione simultanea di tre o più suoni. Nella pratica tonale si costruisce quasi sempre sovrapponendo terze. Un accordo di tre note è detto \textbf{triade}, uno di quattro note \textbf{quadriade} (o accordo di settima).
\end{definizione}

La costruzione per terze riflette la struttura della serie armonica (Sezione \ref{sec:natura_suono}) e produce suoni che l'orecchio sente come un blocco unico. Nel discorso tonale ogni accordo ha una funzione propria e, se instabile, tende a risolvere su un altro accordo.

\section{Le quattro specie di triade}
\label{sec:costruzione_triade}

\begin{definizione}[La triade]
	\label{def:triade}
	Una \textbf{triade} è un accordo di tre note costruito sovrapponendo due terze a partire da una nota base, detta \textbf{fondamentale}. Le altre due note si chiamano \textbf{terza} e \textbf{quinta}, secondo l'intervallo che formano con la fondamentale.
\end{definizione}

I due tipi di terza, maggiore e minore, combinati danno le quattro specie di triade.

\begin{definizione}[Specie di triadi]
	\label{def:specie_triadi}
	\begin{itemize}
		\item \textbf{Triade maggiore:} terza maggiore e terza minore. Nasce direttamente dalla serie armonica.
		\item \textbf{Triade minore:} terza minore e terza maggiore. È stabile, con un colore più scuro.
		\item \textbf{Triade diminuita:} due terze minori. Contiene il tritono (quinta diminuita), dissonante.
		\item \textbf{Triade eccedente (o aumentata):} due terze maggiori. La quinta aumentata la rende ambigua e instabile.
	\end{itemize}
\end{definizione}

Solo la maggiore e la minore sono \textbf{consonanti} e possono servire da punto di riposo. La diminuita e l'eccedente sono \textbf{dissonanti} e richiedono risoluzione. Perché i rapporti di frequenza spieghino questa differenza è detto nell'Approfondimento \ref{phys:rapporti_frequenza_triadi}.

\begin{esempio}[Le quattro specie principali di triadi]
	\label{ex:quattro_triadi}
	\musica{es_cap02_01}
\end{esempio}

\section{Le triadi sui gradi della scala}
\label{sec:triadi_sui_gradi}

Ogni grado della scala può essere la fondamentale di una triade, e la sua specie dipende dalla struttura della scala. Nel modo maggiore:

\begin{definizione}[Triadi diatoniche nel modo maggiore]
	\label{def:triadi_diatoniche_maggiore}
	\begin{itemize}
		\item \textbf{I, IV, V:} triadi maggiori;
		\item \textbf{II, III, VI:} triadi minori;
		\item \textbf{VII:} triade diminuita.
	\end{itemize}
\end{definizione}

\begin{esempio}[Le sette triadi diatoniche di Do maggiore]
	\label{ex:triadi_diatoniche_do}
	\musica{es_cap02_02}
\end{esempio}

Nell'Esempio \ref{ex:triadi_diatoniche_do} i gradi sono indicati con numeri romani maiuscoli, e sotto è specificata la specie dell'accordo (M maggiore, m minore, dim diminuito).

Nel modo minore le triadi si costruiscono sulla scala minore armonica, quella con la sensibile alzata (Sezione \ref{sec:modi}), perché è la sensibile a dare forza alla dominante.

\begin{definizione}[Triadi diatoniche nel modo minore]
	\label{def:triadi_diatoniche_minore}
	Sulla scala minore armonica:
	\begin{itemize}
		\item \textbf{I, IV:} triadi minori;
		\item \textbf{V, VI:} triadi maggiori;
		\item \textbf{II, VII:} triadi diminuite;
		\item \textbf{III:} triade eccedente.
	\end{itemize}
	In La minore: La-Do-Mi, Si-Re-Fa, Do-Mi-Sol$\sharp$, Re-Fa-La, Mi-Sol$\sharp$-Si, Fa-La-Do, Sol$\sharp$-Si-Re.
\end{definizione}
